% Reproducibility appendix. Every digest, version and command below is checked against the
% repository by tests/test_paper_reproducibility.py. No model result in this section.
\section{Reproducibility}
\label{sec:repro}

Code, data generator, verifier, protocol, frozen evaluation code, reports and per-sample result
flags are at \url{https://github.com/efeerdogmus0/trace-graph-repair}.

\paragraph{Data.} Each build is regenerated from a YAML config with
\texttt{python -m trace\_vr generate --config configs/data\_<name>.yaml}
(generator \texttt{trace-vr-gen-1}). The manifest records a content digest over every sample record (graph,
question, answer and edit label), and a separate digest of the image bytes. The content digest
is the reproducibility contract: the launchers regenerate the data on the GPU machine and refuse
to train or evaluate unless it matches the committed manifest. Image bytes are not part of the
contract, because PNG encoding differs between machines; the image digest is recorded per run.

\begin{table}[h]
\centering\small
\begin{tabular}{@{}lll@{}}
\toprule
Build & Config & Content digest (SHA-256, first 16 hex) \\
\midrule
smoke & \texttt{data\_smoke.yaml} & \texttt{c26d05b470bdd5ea} \\
pilot & \texttt{data\_pilot.yaml} & \texttt{ab499a7351a81aaa} \\
pilot\_test & \texttt{data\_pilot\_test.yaml} & \texttt{f7c80f9383d273c1} \\
pilot\_diag & \texttt{data\_pilot\_diag.yaml} & \texttt{bafc07abe3c4743e} \\
\bottomrule
\end{tabular}
\caption{Committed builds. Full digests are in \texttt{manifests/<name>.json}.}
\label{tab:digests}
\end{table}

\paragraph{Training and development evaluation.} One command trains an arm and evaluates it on
development roots:
\texttt{bash scripts/colab/job.sh <job> <arm> <roots> <mode>}, with $3$ epochs, gradient
accumulation $4$, seed $3407$, a checkpoint every $256$ steps and at most $96$ new tokens at
evaluation. Rerunning the same command resumes from the last checkpoint and skips finished
stages. On Kaggle the same script runs inside a private kernel built by
\texttt{scripts/kaggle/make\_kernel.py}, which embeds a byte-reproducible tarball of the
repository and pins \texttt{transformers==5.22.2} and \texttt{peft==0.21.2}.
The six final kernels that print their environment all ran on the same Kaggle image with
\texttt{torch 2.11.0+cu130}, CUDA \texttt{13.0}, Python \texttt{3.13.16} and Tesla T4 GPUs, of which the
kernel uses one; the recorded lines are collected in \texttt{reports/final/kernel\_env.json}. The final
kernels built before this probe was added, and the clean-clone rerun, did not record it.

\paragraph{Locked test evaluation.} The test builds cannot be read by any script unless the
environment variable \texttt{TRACE\_UNLOCK\_TEST} is set to \texttt{final-evaluation}.
\texttt{scripts/colab/final\_eval.sh <job> <run> --final-evaluation} sets it, regenerates the
three builds it needs, checks their content digests and evaluates the six locked splits.

\paragraph{Scoring outside the GPU machine.} Every evaluation also writes per-sample
correctness as a bit string; intervals, paired differences and slices are recomputed locally
from these files, and the collector refuses a run whose local re-score differs from the
summary printed on the GPU machine.

% Generated by scripts/compute_table.py (rule compute-table-v1) from reports/budget.json; do not edit by hand.
\newcommand{\computeCash}{0}
\newcommand{\computeMeasuredHours}{74.33}
\newcommand{\computeMeasuredRuns}{47}
\newcommand{\computeLogLines}{55}
\newcommand{\computeUnmeasuredRuns}{5}
\newcommand{\computeFailedRuns}{2}
\newcommand{\computeFailedMinutes}{6.5}
\newcommand{\computeDeviceUnrecorded}{14}
\newcommand{\computeFinalHours}{41.15}
\newcommand{\computeFinalKernels}{13}
\newcommand{\computeTable}{%
\begin{tabular}{@{}lrrr@{}}
\toprule
Purpose & lines timed & GPU hours & lines not timed \\
\midrule
Smoke and overfit checks & 4 & 0.40 & 0 \\
One-epoch pilot (arms A, B, D) & 5 & 0.97 & 0 \\
Three-epoch dev pilot (all arms, seeds, oracle) & 18 & 16.42 & 8 \\
Frozen-base prompt scan (dev) & 1 & 0.64 & 0 \\
External evaluation (VGCure) & 5 & 13.83 & 0 \\
Locked-test evaluation kernels & 13 & 41.15 & 0 \\
Clean-clone GPU rerun (dev, arm D) & 1 & 0.93 & 0 \\
\midrule
Google Colab & 14 & 4.10 & 8 \\
Kaggle & 33 & 70.24 & 0 \\
\midrule
Total & 47 & 74.33 & 8 \\
\bottomrule
\end{tabular}}

\paragraph{Compute.} No money was spent (cash cost $\computeCash$); all GPU time came from Kaggle's
free weekly quota and Google Colab. The budget log records, for each run, the training plus evaluation
time printed by the job itself. Summed over the \computeMeasuredRuns{} timed log lines this is
\computeMeasuredHours{} GPU hours (Table~\ref{tab:compute}), of which \computeFinalHours{} hours are the
\computeFinalKernels{} locked-test evaluation kernels collected so far and \computeFailedMinutes{} minutes
are \computeFailedRuns{} runs that failed or crashed. Every recorded device is an NVIDIA T4;
\computeDeviceUnrecorded{} Kaggle lines did not record the device. The total is a lower bound: setup,
model download and data generation on the GPU machine are not timed, and the
\computeUnmeasuredRuns{} Colab runs lost with their virtual machine before read-back count as zero.
The table and every number in this paragraph are generated by \texttt{scripts/compute\_table.py} from
\texttt{reports/budget.json}, which \texttt{scripts/budget\_report.py} builds from \texttt{runs/budget.jsonl}.

\begin{table}[h]
\centering\small
\computeTable
\caption{Measured GPU time by purpose and by platform, from the budget log. Lines not timed are runs
lost before read-back and placeholder lines written at launch; both count as zero.}
\label{tab:compute}
\end{table}

\paragraph{Licences.} The source code, tests, notebooks and demo are released under the
MIT licence (\texttt{LICENSE}). The generated benchmark data and images,
their manifests, the reports, figures and paper text, and the trained LoRA adapter weights are released under
CC BY 4.0 (\texttt{LICENSE-DATA.md}). Neither licence covers the base model, whose weights are not
redistributed and keep their own licence (Apache-2.0 per its model card), the DejaVu fonts used to draw labels
(their own licence, shipped with them), or the external benchmark of Section~\ref{sec:external}, of which only
identifiers, counts and hashes are kept. \texttt{THIRD\_PARTY.md} lists every third-party component and the
licence it declares.
